\chapter{Feature, scenari e storie}
\label{cap:feature}

\begin{inbreve}
In un prodotto i requisiti non sono fissati da un cliente: bisogna scoprire che cosa serve agli utenti. Il capitolo
presenta gli strumenti per farlo, in ordine: le \emph{personas}, gli \emph{scenari}, le \emph{user stories}. Da
essi si ricava l'elenco delle \emph{feature} del prodotto. Il capitolo spiega anche come sceglierle bene ed evitare
che il loro numero cresca senza controllo. Le user stories alimentano il product backlog del capitolo~\ref{cap:agile}.
\end{inbreve}

\section{Perché servono le feature}
\label{sec:feature-intro}

Il design di un prodotto software è guidato da quattro fattori:
\begin{itemize}
  \item l'ispirazione;
  \item bisogni di aziende o consumatori non soddisfatti dai prodotti esistenti;
  \item l'insoddisfazione per i prodotti esistenti;
  \item cambiamenti tecnici che rendono possibili nuovi tipi di prodotto.
\end{itemize}

Il \ingl{product-based software engineering} richiede meno documentazione dei requisiti di quello su progetto
(\rif{sec:prodotti-progetti}), perché i requisiti non li fissa un cliente e cambiano. Si identificano invece le
\termine{feature} del prodotto, cioè frammenti di funzionalità.

Per farlo bisogna capire gli utenti potenziali, con interviste, sondaggi e consultazioni informali. Attenzione:
in un'azienda i manager comprano il prodotto, ma i dipendenti che lo usano possono non volerlo usare.

Tre strumenti aiutano a individuare le feature (\figrif{fig:feature-relazioni}): le rappresentazioni degli utenti
(\termine{personas}) e le descrizioni in linguaggio naturale (\termine{scenari} e \termine{storie}).

\begin{figure}[ht]
\centering
\begin{tikzpicture}[blocco/.append style={minimum width=2.6cm}]
\node[blocco] (p) at (0,3.2) {Personas};
\node[blocco] (sc) at (0,1.6) {Scenarios};
\node[blocco] (st) at (4.6,1.6) {Stories};
\node[blocco] (f) at (2.3,0) {Features};
\draw[freccia] (p) -- node[right, font=\footnotesize] {inspire} (sc);
\draw[freccia] (sc) -- node[above, font=\footnotesize] {are-developed-into} (st);
\draw[freccia] (sc) -- node[below left, font=\footnotesize] {inspire} (f);
\draw[freccia] (st) -- node[below right, font=\footnotesize] {define} (f);
\draw[freccia, dashed] (p) -- (st);
\end{tikzpicture}
\caption{Dagli utenti alle feature. Le personas ispirano gli scenari, gli scenari si sviluppano in storie e
ispirano le feature, le storie definiscono le feature. La freccia tratteggiata indica che si può passare
direttamente dalle personas alle storie (\rif{sec:feature-scenari-storie}).}
\label{fig:feature-relazioni}
\end{figure}

\section{Personas}
\label{sec:feature-personas}

Per progettare feature utili bisogna capire chi sono gli utenti, perché contano il loro background, le loro competenze e la
loro esperienza.

\begin{definizione}[Persona]
Una \termine{persona} è un tipo (immaginato) di utente del prodotto. Per esempio, per l'agenda di uno studio
dentistico: il dentista, la segretaria, il paziente.
\end{definizione}

Di solito bastano poche personas (1--2, al massimo 5) per identificare le feature chiave.

\subsection{Che cosa contiene la descrizione di una persona}

\begin{description}[leftmargin=0pt, style=unboxed]
  \item[Personalization.] Un nome e qualche informazione sulla situazione personale. A volte è utile una fotografia
    adatta (\rif{sec:feature-foto}).
  \item[Job-related.] Se il prodotto è destinato alle aziende, qualcosa sul lavoro della persona e, se serve, su ciò
    che comporta.
  \item[Education.] La formazione e il livello di competenze tecniche e di esperienza. È importante soprattutto per
    il design dell'interfaccia.
  \item[Relevance.] Se possibile, perché la persona potrebbe essere interessata al prodotto e che cosa potrebbe volerne
    fare.
\end{description}

\begin{esempio}[Emma, a history teacher]
Emma ha 41 anni e insegna storia in una scuola superiore di Edimburgo a studenti di 12--18 anni. È nata a Cardiff,
dove i genitori erano insegnanti. Dopo la laurea in storia si è trasferita a Edimburgo e ha seguito la formazione
da insegnante. Ha due figli di 6 e 8 anni. Torna a casa presto per vederli, quindi prepara le lezioni, svolge le pratiche amministrative e corregge i compiti da casa.

Usa i social media e le applicazioni di produttività abituali per preparare le lezioni. Non è particolarmente
interessata alle tecnologie digitali. Odia l'ambiente di apprendimento virtuale della scuola e lo evita quando può.
Crede che l'insegnamento in presenza sia il più efficace. Potrebbe usare il sistema iLearn per l'amministrazione e
per accedere a film e documenti storici, ma non è interessata a un approccio misto digitale/in presenza.

Le conclusioni sono due: \emph{non ha un background tecnico} e \emph{vuole un prodotto che supporti
l'amministrazione}.
\end{esempio}

\begin{idea}
Le personas permettono agli sviluppatori di ``mettersi nei panni degli utenti''. Quando studiare gli utenti non è
possibile (per esempio per alcuni prodotti nuovi) si usano le \termine{proto-personas}, costruite senza
dati diretti sugli utenti.
\end{idea}

\subsection{Le fotografie nelle personas}
\label{sec:feature-foto}

L'uso di una fotografia è discusso. Le osservazioni sono quattro:
\begin{itemize}
  \item a favore: alcuni studi suggeriscono che una fotografia aiuta i team a usare le personas con più efficacia; i
    profili con una foto sorridente aumentano la disponibilità a usarle (J.~Salminen e altri);
  \item contro: le personas non dovrebbero riguardare l'aspetto delle persone, ma ciò che fanno (S.~Cable);
  \item contro: personas molto dettagliate spingono il team a supporre che i dati demografici guidino le motivazioni
    (S.~Wachter-Boettcher);
  \item il genere ha un'influenza mediatrice sugli attributi percepiti di uomini e donne (F.~Deutsch).
\end{itemize}

\section{Scenari}
\label{sec:feature-scenari}

\begin{definizione}[Scenario]
Uno \termine{scenario} è una narrazione che descrive una situazione in cui un utente usa le feature del prodotto per
fare qualcosa che vuole fare. Serve a scoprire le feature.
\end{definizione}

La descrizione di uno scenario contiene sei elementi:
\begin{itemize}
  \item il nome dello scenario;
  \item le personas (attori) coinvolte;
  \item l'obiettivo generale;
  \item il problema che il sistema esistente non può affrontare;
  \item che cosa comporta raggiungere l'obiettivo;
  \item i modi possibili per affrontare il problema.
\end{itemize}

\begin{esempio}[Scenario ``Fishing in Ullapool'']
Jack è un insegnante di scuola primaria a Ullapool e ha deciso che un progetto di classe riguardi l'industria della
pesca della zona. Gli studenti raccolgono ricordi dai parenti, usano archivi di giornali e fotografie e le condividono
in un wiki del sistema iLearn. Jack ha bisogno di un sito per condividere foto, perché gli studenti devono commentare
quelle degli altri. Deve poter moderare i post con foto prima che siano condivisi, per i problemi di copyright e
privacy. Chiede consiglio al gruppo di insegnanti di scuola primaria. Due colleghi gli suggeriscono KidsTakePics.
Poiché il sito non è integrato con l'autenticazione di iLearn, Jack crea un account da insegnante e uno di classe.
Aggiunge KidsTakePics ai servizi visibili agli studenti, così che al login possano subito caricare foto da telefono e
dai computer di classe.
\end{esempio}

Gli scenari narrativi, di alto livello, facilitano la comunicazione e stimolano la creatività nel design. Non sono però
specifiche: mancano di dettaglio e possono essere incompleti.

\subsection*{Come si scrivono}

\begin{itemize}
  \item Più scenari (per esempio 3--4) per ogni persona, che coprono le sue responsabilità principali.
  \item Sono scritti dal punto di vista dell'utente.
  \item Ogni membro del team ne crea alcuni individualmente. Poi li discute con il resto del team e, se possibile, con
    gli utenti.
\end{itemize}

\section{User stories}
\label{sec:feature-storie}

\begin{definizione}[User story]
Gli scenari sono storie di uso del prodotto \emph{di alto livello}. Le \termine{user stories} sono narrazioni più
fini, che seguono il modello:

\smallskip
\begin{tabular}{@{}ll@{}}
\textbf{As a} & $\langle$ruolo$\rangle$ \\
\textbf{I want to} & $\langle$fare qualcosa$\rangle$ \\
\textbf{So that} & $\langle$motivo / valore$\rangle$ \\
\end{tabular}
\end{definizione}

Le storie realizzano tre principi del manifesto agile (\rif{sec:agile-manifesto}): consegnare presto software di
valore (1), consegnare frequentemente (3), cercare la semplicità (10). Per questo:
\begin{itemize}
  \item il lavoro è organizzato in unità che rappresentano valore per il cliente, distinguendo ciò che vale di più;
  \item le unità di lavoro sono visualizzate con storie semplici, scritte dal punto di vista dell'utente;
  \item storie di buona qualità rendono più efficiente il resto dello sviluppo e creano valore di business per gli utenti;
  \item il software si costruisce dalla prospettiva dell'utente.
\end{itemize}

\subsection*{Suddividere le storie}

Gli sviluppatori tendono a vedere l'applicazione come uno \emph{stack}. L'utente la vede come \termine{fette di
funzionalità} di valore. Quindi:
\begin{itemize}
  \item si costruisce il software in modo incrementale, a fette di funzionalità;
  \item una buona storia richiede 1--2 giorni di lavoro;
  \item le storie lunghe (\termine{epics}) si spezzano in storie più brevi, che si possono prioritizzare.
\end{itemize}

\begin{esempio}[Spezzare una epic\aggiunta]
La storia del Tic Tac Toe \emph{``As a Player I want to play against the computer so that I can learn to play
alone''} richiede molto più di 2 giorni: bisogna progettare il comportamento del computer. Si spezza in tre fette,
ciascuna utilizzabile da sola:
\begin{enumerate}
  \item il computer gioca una mossa valida qualsiasi;
  \item il computer blocca la mossa con cui l'avversario vincerebbe;
  \item il computer vince quando può farlo.
\end{enumerate}
Dopo la prima fetta il giocatore può già giocare da solo, anche se male: c'è valore consegnabile. Le altre due si
ordinano per priorità e si consegnano negli sprint successivi.
\end{esempio}

\subsection*{Storie e pianificazione}

Il product backlog di Scrum (\rif{sec:agile-backlog}) è spesso un insieme di user stories. Le storie sono associate a
priorità (e talvolta a una stima dello sforzo) e ordinate in base a esse. Questo ordinamento si chiama
\ingl{requirements triage}.

\section{Storie e scenari}
\label{sec:feature-scenari-storie}

Tutte le funzionalità descritte in uno scenario si possono esprimere come user stories. Gli scenari restano utili perché:
\begin{itemize}
  \item si leggono in modo più naturale;
  \item rendono più facile capire le storie;
  \item danno più contesto.
\end{itemize}
D'altra parte si può passare direttamente dalle personas alle storie, saltando gli scenari (\figrif{fig:feature-relazioni}).

\begin{esempio}[Tic Tac Toe]
\emph{Regole}: griglia 3$\times$3; due giocatori (X e O); a turno mettono il proprio segno in una casella vuota; vince
chi per primo ottiene 3 segni in riga, colonna o diagonale; quando le 9 caselle sono piene la partita finisce.

\emph{Persona (Player)}: Pierfrancesco, 9 anni, scolaro delle elementari a Pisa. Vive con i genitori e una sorella
maggiore e gli piacciono la spiaggia, i videogiochi e il calcio. Sul treno per Firenze, dove abitano i nonni, gioca a
giochi da tavolo con il padre. Sa usare computer e smartphone ma non vuole approfondire la tecnologia: per lui usarla
serve a divertirsi. Gli piacerebbe giocare ai giochi da tavolo sul telefono o sul PC. (Persona generata dal corso
dell'anno precedente.)

\emph{User stories} (tutte con ruolo \emph{Player}):

\smallskip
\small
\begin{tabularx}{\linewidth}{>{\raggedright\arraybackslash}p{4.2cm}X}
\toprule
I want to & So that \\
\midrule
have a timer & the game is even shorter \\
select my next valid move & I can play the game \\
select my symbol & I can customise my games \\
play against the computer & I can (learn how to) play alone \\
know when I win/fail/tie & I can celebrate, be sad, or play again \\
ask for a rematch & I can eventually win \\
surrender & I can stop suffering from losing \\
see an updated board & I can decide on my next move \\
see whose turn it is & I know if I should play a move \\
see the winning combo & I can better see it \\
start a new game & I can play with others on same device \\
\bottomrule
\end{tabularx}
\end{esempio}

\section{Identificare le feature}
\label{sec:feature-identificazione}

L'obiettivo è ottenere l'elenco delle feature che definiscono il prodotto. Una buona feature ha tre proprietà.

\begin{definizione}[Proprietà delle feature]
\begin{itemize}
  \item \textbf{Independence}: non dipende da come sono implementate le altre feature e non è influenzata dall'ordine
    di attivazione delle altre.
  \item \textbf{Coherence}: è legata a una sola funzionalità. Non fa più di una cosa e non ha mai effetti collaterali.
  \item \textbf{Relevance}: riflette il modo in cui gli utenti svolgono normalmente un compito. Non offre
    funzionalità oscure, richieste di rado.
\end{itemize}
\end{definizione}

\subsection*{Fonti di conoscenza}

Il design delle feature si basa su quattro fonti di conoscenza.

\begin{description}[leftmargin=0pt, style=unboxed]
  \item[User knowledge.] Scenari e storie informano il team su che cosa vogliono gli utenti e su come potrebbero usare le
    feature.
  \item[Product knowledge.] L'esperienza con prodotti esistenti, o una ricerca su ciò che fanno. A volte le feature
    devono replicare quelle di questi prodotti, perché offrono funzionalità fondamentali sempre richieste.
  \item[Domain knowledge.] La conoscenza dell'area di lavoro (per esempio la finanza o la prenotazione di eventi) che
    il prodotto vuole supportare. Capire il dominio permette di pensare a modi nuovi di aiutare gli utenti.
  \item[Technology knowledge.] I nuovi prodotti nascono spesso per sfruttare sviluppi tecnologici successivi
    al lancio dei concorrenti. Conoscere la tecnologia più recente permette di progettare feature che la sfruttano.
\end{description}

\subsection*{Fattori nella progettazione dell'insieme di feature}

Nella scelta dell'insieme di feature ci sono tre coppie di fattori in tensione (\figrif{fig:feature-fattori}).

\begin{figure}[ht]
\centering
\begin{tikzpicture}
\node[font=\small\bfseries\sffamily] at (0,2.2) {Feature set design factors};
\node[font=\small] (s) at (-4,1.3) {Simplicity};
\node[font=\small] (fu) at (4,1.3) {Functionality};
\node[font=\small] (co) at (-4,0) {Control};
\node[font=\small] (au) at (4,0) {Automation};
\node[font=\small] (fa) at (-4,-1.3) {Familiarity};
\node[font=\small] (no) at (4,-1.3) {Novelty};
\draw[<->, thick, rosso] (s) -- (fu);
\draw[<->, thick, rosso] (co) -- (au);
\draw[<->, thick, rosso] (fa) -- (no);
\end{tikzpicture}
\caption{Fattori nel design dell'insieme di feature. Ogni freccia lega due fattori in tensione: aumentare l'uno
tende a ridurre l'altro (simplicity contro functionality, control contro automation, familiarity contro novelty).}
\label{fig:feature-fattori}
\end{figure}

\section{Feature creep e derivazione delle feature}
\label{sec:feature-creep}

Il numero di feature cresce quando si immaginano nuovi utenti potenziali. Questo fenomeno si chiama
\termine{feature creep}. Le cause sono tre:
\begin{itemize}
  \item i dirigenti di marketing vogliono soddisfare tutte le richieste degli utenti;
  \item la pressione del marketing a includere le feature dei concorrenti;
  \item il desiderio di supportare sia utenti esperti sia inesperti.
\end{itemize}

Per evitarlo, per ogni feature proposta ci si pone quattro domande:
\begin{itemize}
  \item aggiunge davvero qualcosa di nuovo, o è solo un modo alternativo di fare ciò che è già supportato?
  \item è probabile che sia importante e usata dalla maggior parte degli utenti?
  \item si può implementare estendendo una feature esistente, invece di aggiungerne un'altra al sistema?
  \item offre una funzionalità generale o è una feature molto specifica?
\end{itemize}

\subsection*{Derivazione e descrizione delle feature}

Il team di sviluppo si riunisce per discutere scenari e storie ed estrarne l'elenco delle descrizioni delle feature.
Ogni descrizione indica (\figrif{fig:feature-descrizione}):
\begin{itemize}
  \item il \emph{nome} della feature;
  \item l'\emph{input}, cioè i dati dall'utente e da altre fonti;
  \item l'\emph{activation}, cioè come l'utente attiva la feature;
  \item l'\emph{action}, cioè come sono elaborati i dati di input;
  \item l'\emph{output}, cioè ciò che è restituito all'utente e al sistema.
\end{itemize}
Si inizia poi a costruire prototipi (\rif{sec:prodotti-prototipo}) per dimostrare le feature, prima di tutto quelle
\emph{nuove} e quelle \emph{critiche}.

\begin{figure}[ht]
\centering
\begin{tikzpicture}[blocco/.append style={minimum width=3.4cm}]
\node[blocco] (i) at (-3.2,2.2) {Input\\dall'utente e da\\altre fonti};
\node[blocco] (a) at (3.2,2.2) {Activation\\come l'utente attiva\\la feature};
\node[blocco] (az) at (0,0.6) {Action\\come sono elaborati\\i dati di input};
\node[blocco] (o) at (0,-1.2) {Output\\all'utente e\\al sistema};
\node[font=\small\bfseries\sffamily] at (0,2.8) {Feature name};
\draw[freccia] (i) -- (az);
\draw[freccia] (a) -- (az);
\draw[freccia] (az) -- (o);
\end{tikzpicture}
\caption{Struttura della descrizione di una feature: nome, input, activation, action e output.}
\label{fig:feature-descrizione}
\end{figure}

\begin{puntichiave}
\begin{itemize}
  \item Un prodotto ha meno documentazione dei requisiti: si identificano le \emph{feature}, frammenti di funzionalità.
  \item Personas $\rightarrow$ scenari $\rightarrow$ storie $\rightarrow$ feature. Si può saltare gli scenari.
  \item Persona: tipo di utente immaginato, 1--2 (al massimo 5) personas, con personalization, job-related, education,
    relevance.
  \item Scenario: narrazione di alto livello, non una specifica; 3--4 per persona, dal punto di vista dell'utente.
  \item User story: As a $\langle$ruolo$\rangle$ I want to $\langle$azione$\rangle$ So that $\langle$motivo$\rangle$.
    Una storia richiede 1--2 giorni; le epics si spezzano.
  \item Una feature deve essere independent, coherent, relevant.
  \item Il feature creep si evita con quattro domande su novità, importanza, estendibilità e generalità.
\end{itemize}
\end{puntichiave}

\begin{domande}
\begin{enumerate}
  \item Perché un prodotto richiede meno documentazione dei requisiti di un progetto?
  \item Che cos'è una persona e che cosa contiene la sua descrizione? Quante ne servono?
  \item In che cosa differiscono scenari e user stories? Perché usare ancora gli scenari?
  \item Qual è il modello di una user story? Quanto deve durare e che cos'è una epic?
  \item Quali sono le tre proprietà di una buona feature?
  \item Che cos'è il feature creep, che cause ha e come si evita?
  \item Che cosa contiene la descrizione di una feature?
\end{enumerate}
\end{domande}
